\section{What was verified by machine}\label{sec:verification}

Apart from the statements transcribed in \cref{sec:statements}, the description of the
elaborators in \cref{sec:scan}, the account of the reviews in \cref{sec:checker}, and what
is said of how Lean and \lean{lean4lean} work, which rests on the reviews and on our reading
of their source, every result in this section was computed by a program: on the archive, or,
where we say so, on Mathlib and on test modules. The programs and their output are public
(\cref{app:reproduce}).

\subsection{The archive}
The Lean source is distributed as fourteen parts which concatenate to a single file
\lean{lean.tar.gz} of \archivebytes{} bytes with SHA-256
\[
\text{\small\ttfamily \archiveshaA}\;\text{\small\ttfamily \archiveshaB}.
\]
It unpacks to \filecount{} \lean{.lean} files (\linecount{} lines) and \otherfiles{} build
files. The build files pin Lean \leanversion{} and Mathlib at commit
{\small\ttfamily \mathlibrev}.

\subsection{Static inspection}\label{sec:scan}
We scanned the text of every file, comments included, for whole words, also where a word
follows a dot. No file contains the word \lean{sorry} or \lean{admit}, an \lean{axiom}
declaration, or any of \lean{native_decide}, \lean{implemented_by}, \lean{extern},
\lean{unsafe}, \lean{opaque}, \lean{csimp}; nor any of the names we know by which a
declaration can be put into an environment without the kernel's check or by which the
compiler is trusted, which the script lists (among them \lean{skipKernelTC},
\lean{addDeclWithoutChecking}, \lean{doCheck} and \lean{ofReduceBool}). No file begins with
the word \lean{module}, so the archive does not use Lean's module system
(\cref{sec:checker}). Every import names a file of the archive or a module of Mathlib, of its
dependency Batteries or of Lean, and the build shows that those modules exist; the import
graph is acyclic; every file is in the import closure of \lean{Main.lean}.

Counted by keyword, the text has \scanmacros{} macros, which abbreviate syntax, and
\scanelabs{} elaborators, programs that run while a file is compiled, in \scanelabfiles{}
files. We read them: they are two tactics, a command that copies declarations under new
names, and a probe. Each produces its result through an ordinary declaration, so the kernel
checks what they produce; the copying command is the one place where the text calls Lean's
function for adding a declaration, the one that checks. In \scanruncmdfiles{} small files a
command is run while the file is compiled; it adds a name to the deny list of Lean's facility
for suggesting library lemmas, and declares nothing. The elaborators have
\scanpartial{} helper functions declared \lean{partial}; they are programs, not statements,
and they are outside the dependency cone of \cref{sec:cone}. No file contains \lean{#eval}.
A scan of the text can miss a defect, and nothing below rests on it: the replay of
\cref{sec:replay} examines what the build produced, however it was produced.

\subsection{The statements}\label{sec:statements}
\lean{Main.lean} is \mainlines{} lines long and proves five theorems.
\begin{statement}[\lean{riemannHypothesis}]\label{st:rh}
For every $s\in\C$ with $\zeta(s)=0$, if $s$ is not of the form $-2(n+1)$ with $n\in\N$ and
$s\neq 1$, then $\re s=\tfrac12$.
\end{statement}
The statement is the proposition \lean{RiemannHypothesis} of Mathlib, not a definition of the
archive: in the compiled environment that name is declared in the module
\lean{Mathlib.NumberTheory.LSeries.RiemannZeta}, as is \lean{riemannZeta}.
\begin{statement}[\lean{LFunction_ne_zero_of_half_lt_re}]\label{st:grh}
For every modulus $N\ge 1$, every Dirichlet character $\chi$ modulo $N$ and every $s\in\C$ with
$\re s>\tfrac12$, if $\chi\neq 1$ or $s\neq1$ then $L(s,\chi)\neq0$.
\end{statement}
\begin{statement}[\lean{LFunction_eq_zero_re}]\label{st:strip}
For every modulus $N\ge1$, every Dirichlet character $\chi$ modulo $N$ and every $s$ with
$0<\re s<1$ and $L(s,\chi)=0$, one has $\re s=\tfrac12$.
\end{statement}
The remaining two are \lean{riemannZeta_eq_zero_re}, which is \cref{st:rh} again, and
\lean{LFunction_eq_zero_re_of_isPrimitive}, which is \cref{st:strip} for primitive characters.
Here $L(s,\chi)$ is Mathlib's \lean{DirichletCharacter.LFunction}, the meromorphic continuation.
The characters are not assumed primitive; restricting to the open strip in \cref{st:strip}
excludes the zeros on $\re s=0$ that imprimitive characters have, so the statements are
correctly posed.

``GRH'' in this paper always means \cref{st:grh,st:strip}: Dirichlet $L$-functions. The note
says the case of Dedekind zeta functions is expected and not done.

\subsection{The build}\label{sec:build}
\ifsecondrun
We built the unmodified archive twice, on two rented machines. The first build took
Mathlib's compiled files from the community's build cache for the pinned commit and ran
\lean{lake build Main} for \firstbuildhours{} hours on \firstbuildthreads{} threads
(\firstbuildcpu); its log was not kept, and these figures are from a summary of it made by
hand. For the second, Mathlib and the packages it depends on were first compiled from
source, \mathlibmodules{} modules of Mathlib in \mathlibminutes{} minutes, the download of
the toolchain included, so that no compiled file came from a cache; \lean{lake build Main}
then ran for \buildhours{} hours on \buildthreads{} threads (\buildcpu), while most of the
tests of \cref{sec:checker} on Mathlib and on the test modules ran beside it. Each build ended with
the line \lean{Build completed successfully}, after \buildjobs{} jobs in the second, and
with exit code~$0$. In each, all \filecount{} modules of the archive were built and none
failed. On the second machine the installed toolchain was compared, file by file, with the
release archive of Lean \leanversion{} published on GitHub: all \toolchainfiles{} files are
identical. The first machine's toolchain was the same release, installed in the same way;
it was not compared.

The two builds give the same names for the archive. The list of its declarations inside the
cone of \cref{sec:cone}, each with its kind, the list of those outside it, and the list of
its modules that meet the cone are the same in both, byte for byte, and so are the counts of
the two replays for every kind of declaration but one. They differ in Mathlib. The compiled
files of the first build were not kept, so the cache was fetched again on the second machine
and compared with the build from source: for every constant its kind, its module and its
name, and Lean's hash of its statement and of its value. Compiled from source, Mathlib has
\mathlibonlysource{} theorems that the cached build lacks and lacks \mathlibonlycache{} that
it has. All of these are private auxiliary lemmas that the tactics \lean{simp} and
\lean{abel} produce inside \mathlibdiffproofs{} proofs; no name that a person wrote is among
them. \mathlibhashnote{} The second build thus declares \mathlibextra{} theorems more than
the first, which is the difference between the two replays, and its cone is larger by
\mathlibextraconewords. What follows refers to the second build.
\else
We ran \lean{lake build Main} on the unmodified archive, on a rented machine with
\buildthreads{} threads (\buildcpu), using the Mathlib build cache for the pinned commit. It
ended after \buildhours{} hours with the line \lean{Build completed successfully}, after
\buildjobs{} jobs, and with exit code~$0$. All \filecount{} modules of the archive were built;
none failed. A second build is in progress on another machine, with Mathlib compiled from
source and not taken from the cache; the next version of this paper reports it. The toolchain
installed on that second machine was compared, file by file, with the release archive of Lean
\leanversion{} published on GitHub: all \toolchainfiles{} files are identical. The first
machine's toolchain was the same release, installed in the same way; it was not compared.
\fi

\subsection{Axioms}\label{sec:axioms}
For each of the five theorems, \lean{#print axioms} reports exactly
\[
\text{\lean{propext}},\qquad \text{\lean{Classical.choice}},\qquad \text{\lean{Quot.sound}},
\]
the three axioms on which Mathlib as a whole rests. In particular none depends on
\lean{sorryAx} or on an axiom that trusts the compiler. For a theorem of an imported module
this command does not walk the terms. It prints a list that Lean computed and recorded when
the module that declares the theorem was compiled, from the terms of that module and from the
lists recorded for the modules it imports; it walks the terms only where no list is recorded
(\lean{Lean/Util/CollectAxioms.lean} in the source of Lean \leanversion{}). So the report is
as good as the compiled files it reads, and in the first build those of Mathlib came from the
community's cache.
\ifsecondrun
Two programs of ours compute the axioms again from the terms. The replay of
\cref{sec:replay} does it for each of the five theorems, in the environment that the kernel
built, and finds the same three. The program that computes the cone of \cref{sec:cone} finds
these three and no other axiom among the constants of the cone of the five theorems.
\else
The corrected replay program and the program that computes the cone of \cref{sec:cone} both
compute the axioms from the terms; they are part of the run now under way (\cref{sec:replay}),
and this version has the recorded lists only. The first replay's list does show that the
whole import closure of \lean{Main.lean} declares \replayaxioms{} axioms that are not marked
unsafe, and no more: these three, \lean{sorryAx}, and three that trust the compiler.
\fi

\subsection{The dependency cone}\label{sec:cone}
A build proves that every file compiles; it does not say which declarations the final theorem
uses. From the compiled environment we computed the set of all constants reachable from
\lean{riemannHypothesis} through the type and the value of every declaration reached, the
link between an inductive type and its constructors, and the structure that a projection
belongs to. It has
\conesize{} elements. Of the \projtheorems{} theorems the archive declares,
\coneprojtheorems{} are in it, and the cone meets \coneprojmodules{} of the \projmodules{}
modules that declare anything. These are counts of names: Lean lets an auxiliary theorem be
declared again in each module that needs it, and \repeatedtheorems{} of the archive's theorem
declarations repeat a name declared in another module. When this paper says that a declaration
is, or is not, used by the proof, it means membership in this set. \Cref{sec:offpath} lists
what is outside it.

\subsection{Kernel replay}\label{sec:replay}
Some words of Lean are needed here. The \emph{kernel} is the small program that checks a
declaration. An \emph{inductive type} comes with its \emph{constructors} and with a
\emph{recursor}, its induction principle; the kernel derives these from the declaration of
the type, and it alone provides the four constants of the \emph{quotient type}. An
\emph{opaque constant} is a definition whose value later declarations cannot use. Lean marks
a declaration \emph{unsafe} or \emph{partial} when it is a program and not a statement; no
theorem may use one. A file that begins with the word \lean{module} uses Lean's \emph{module
system}, by which a file shows the files that import it a reduced version of its
declarations. Mathlib and Lean's library use it; the archive does not (\cref{sec:scan}).

A build normally checks each declaration once, in the process that elaborated it; code that a
file runs while it is compiled can add a declaration that the kernel never sees
(\cref{sec:checker}). A replay passes the declarations to the kernel a second time, from the
compiled files, in a process that does nothing else. We replayed the whole import closure of
\lean{Main.lean} into an environment that starts empty: the archive, Mathlib with its
dependencies, and Lean's own library, \replaymodules{} modules that declare
\replayconstants{} distinct names. Of these, \replaydecls{} are neither unsafe nor partial,
and they are what a replay examines, together with the further declarations of names that
more than one module declares.

\ifsecondrun
Each kind is examined in its own way. A theorem, a definition, an opaque constant or an axiom
is passed to the kernel's entry point \lean{Kernel.Environment.addDeclCore} against an
environment that contains only declarations that precede it in a dependency order. An
inductive type is passed with the other types of its mutual block and their constructors,
once for the block, and the kernel generates the constructors and the recursors itself. The
quotient type is one declaration of the kernel's, which provides four constants. When every
declaration has been sent, the environment the kernel has built is compared with what was
read from the compiled files, in both directions: each of the \replaydecls{} constants must
equal, field by field, the constant of the same name in the kernel's environment, and the
kernel's environment must contain nothing else. Terms are compared up to the names and
annotations of bound variables, which the kernel does not consult. So, apart from the unsafe
and partial declarations, the compiled files hold, for the version of each name that Lean's
import keeps and as Lean's reader of compiled files presents them to a file outside Lean's
module system, exactly what the kernel produces from the declarations in them; and every
further version of a repeated name was accepted by the kernel. What a module inside the
module system would receive from the same files is compared separately (\cref{sec:checker});
in this run that covered \replayexported{} constants in \replayexportedmodules{} modules.

\replaystatus

For each of the five theorems the program also reports that it was checked as a theorem and
that nothing it depends on failed, and computes the axioms it depends on from the terms, in
the environment the kernel built. They are \lean{propext}, \lean{Classical.choice} and
\lean{Quot.sound} for each, and the program that computes the cone finds the same three among
the constants of the cone. This matters because \lean{#print axioms}, for a theorem of an
imported module, prints a list that was recorded in the compiled file when it was written; it
does not walk the terms again. For \lean{riemannHypothesis} the replay program counts
\replayconeconstants{} names, following types, values, constructors, the blocks and rules of
the recursors reached, the structures of projections and the names into which numerals and
strings unfold; the cone of \cref{sec:cone}, which follows types, values, the link between an
inductive type and its constructors and the structures of projections, has \conesize{}.
\else
The replay reported in this version is the one made on the first build, with the first form
of our program. A theorem, a definition, an opaque constant or an axiom was passed to the
kernel's entry point \lean{Kernel.Environment.addDeclCore} against an environment that
contained only declarations that precede it in a dependency order. An inductive type was
passed with the other types of its mutual block and with the constructors that the files hold
for them, once for the block; the kernel generates constructors and recursors itself, and the
ones in the files were compared with those. The quotient type is one declaration of the
kernel's, which provides four constants.

\replaystatus

That replay did less than the first two versions of this paper said, as the reviews of the
program showed (\cref{sec:checker}). It did not compare the inductive types and the quotient
constants in the files with the ones the kernel built. It did not check that the kernel's
environment holds nothing but the constants of the files. Where a file listed a
constructor that no file contains, it would have built the type without it. And of each
compiled file of Mathlib and of Lean's library it read only the part that a file outside the
module system reads. So it establishes that every theorem, definition, opaque constant and
axiom in the parts it read is accepted by the kernel, given the inductive types as the kernel
builds them, and that the constructors and recursors in those parts are the kernel's. It
does not exclude a compiled file whose constants differ from what the kernel built in
another way: an inductive type or a quotient constant other than the kernel's; a constructor
or a recursor that the kernel generates and that the files lack, or bind to a constant of
another kind; a constant stored under a name that is not its own. Through the second of
these, the first review found, a theorem that the kernel would reject against the files
could have been accepted. By the second review's reasoning the counts of that run, which add
up exactly, exclude a further quotient constant and a constant stored under another name,
and exclude none of the other shapes. That review proposed an audit of the compiled files
that would settle the rest. The compiled files of the first build were not kept, so the
audit could not be made on them; the corrected program makes it on the second build.

The program has been
corrected: after the last declaration it now compares the environment the kernel built with
what it read from the compiled files, constant by constant and in both directions, it
refuses an inductive type whose listed constructors are not all there, and it compares what
a module inside the module system would receive. The corrected program
has been run on the test modules and on Mathlib (\cref{sec:checker}). Its run over the whole
archive is under way on the second build, and the next version of this paper reports it.
\fi

\subsection{The checker}\label{sec:checker}
\ifsecondrun
The replay was done by a program of ours, of \replayprogramlines{} lines, published with this
paper.
\else
The replay reported above was made by the first form of our program, of \firstprogramlines{}
lines; the corrected program has \replayprogramlines{} lines. Both are published with this
paper.
\fi
It imports once, orders the declarations, and runs the kernel checks of theorems,
definitions, opaque constants and axioms in parallel, each against a snapshot of the
environment as it stood when the declaration's turn came; inductive types and the quotient
type are checked on the main thread. Lean's own checker is \lean{leanchecker}
\cite{lean4checker}.
\ifsecondrun
Its ordinary mode imports the import closure again for every module it
checks. Run as described below, \stocksampleatatime{} processes at a time and beside the
other runs named there, it took \stocksampleminutes{} minutes for \stocksample{} names; at
that rate, checking each of the \filecount{} modules once would take about \stockfulldays{}
days on that machine. Its function for one module, called on one module at a time, took
\stockexactmin{} to \stockexactmax{} seconds on each of the \stockkilledimports{} modules
named below. Its other mode,
\lean{--fresh}, replays one closure in one process, as our program does, but on one thread.
We ran it on the whole import closure of \lean{Main}, on the second machine. It accepted the
closure, ending with exit code~$0$ after \freshmainhours{} hours, \freshmaincpuhours{} of
them processor time, and printing nothing; for about its first two hours three of the runs
below ran beside it: the sample run whose time is quoted, the second implementation on the
closure of the $L$-function, and the checks of single modules. So each of the
\replaydecls{} constants of \cref{sec:replay}, which are those that are neither unsafe nor
partial, each in the version that Lean's import keeps of its name, has been accepted by the
kernel twice from an empty environment, once driven by Lean's own checker and once by our
program. What rests on our programs alone is what Lean's replay does not do: the comparison
of the kernel's environment with the compiled files beyond constructors and recursors, the
check of the further versions of repeated names, the comparison of the files of the module
system, and the axioms computed from the terms.
\else
Its ordinary mode imports the import closure again for every module it checks, \filecount{}
times over; on the first machine an attempt on some two thousand modules at once ended at
the limit of 128~GB of memory. Its other mode, \lean{--fresh}, replays one closure in one
process, as our program does, but on one thread. On Mathlib it used
\stockfreshmathlibcpuminutes{} minutes of processor time and our program
\stockreplaymathlibcpuminutes{}. The first replay of the archive took \replayminutes{}
minutes on \firstreplaythreads{} threads, so in that mode Lean's checker would need a day or
more of one processor for the archive; we intend to make that run on the second build.
\fi
For those we give the grounds for trusting the replay program. The runs behind
\cref{tab:hostile}, and those with \lean{--fresh}
on Mathlib and on the module that declares \lean{RiemannHypothesis}, were made on the second
machine, on its own build of Mathlib from source, while it was building the archive.
Everything else in this subsection was run after the build: the run of
\lean{leanchecker --fresh} on \lean{Main}, the second implementation on the closure of the $L$-function, the
runs on a sample of names of the archive's modules, and the checks of single modules, with
the two tests of the way of calling them. The sample run whose time is quoted ran beside
the run on \lean{Main} and, for its first half hour, beside that run of the second
implementation.

\emph{Review.} The program has been reviewed four times against the source of Lean
\leanversion{}, by further instances of the language model that wrote it; no person has
reviewed it. All four reports are published with this paper. The first two, made
independently of each other, found among other things that the program used for the first
two versions of this
paper compared only constructors and recursors with the kernel's, so that an inductive type or
a quotient constant written directly into a compiled file was counted and never examined; and
that it rebuilt an inductive type from those of its listed constructors that it found,
dropping one that was missing. Neither found a way for this to affect the five theorems, and
Lean's own replay shares the first gap. But a replay with those gaps establishes less than
those versions said. The program was rewritten along the lines the two reviews proposed. The
third review, which had the first two before it, built and ran the rewritten program. It
found one further gap. Mathlib and Lean's library use Lean's module system, in which a module
is compiled to several files, and the program read only the one that a file outside the
module system reads; a false theorem in another would not have been seen. It
also found that two of the program's checks were exercised by no test module. The program
was changed to compare the other files, and test modules were added. The fourth review, of
those changes, found two more routes of the same kind, by which a module inside the module
system could receive a constant that the program had not seen, and found that the test script
could report on files left by an earlier run. The two routes were closed, each with a test
module, and the script was corrected. The reviewer then checked the corrections, building
the program and running the tests, and found no input on which the corrected program reports
success wrongly; that check, with the points it leaves open, is appended to the fourth
report. None of the gaps found by the third
and fourth reviews touches the five theorems, which are stated and checked outside the module
system (\cref{sec:scan}). After that check only the comment at the head of the program was
changed, to correct sentences that the reviewer had named and to add a paragraph on the
memory it uses; its instructions are the ones checked, as a comparison with the program of
version~1.2 of this paper's repository shows. The script that runs the tests was also
changed after the check: it refuses an empty or repeated list of modules and tests the
program's list of axioms. The tests below were made with the program as it now stands,
whose SHA-256 begins {\small\ttfamily \hostileprogramsha}.
\ifsecondrun
So was the replay reported here. The first run, made with the earliest form of the program on
the first build, had given the same counts of modules, definitions, opaque constants, axioms,
inductive types, constructors and recursors, and \mathlibextra{} theorems fewer, the
difference between the two builds of Mathlib (\cref{sec:build}).
\fi

\emph{Modules built to be wrong.} We made \hostilemodules{} small modules
(\cref{tab:hostile}). \hostilecrafted{} are compiled files written directly; every other one
has a source and compiles. \hostilebad{} of them have, in their import closure, a compiled
file that holds
something the kernel never accepted, or whose parts disagree in a way that Lean never writes:
a theorem added without a check, a recursor altered after
the kernel made it, a constant of a kind that only the kernel creates. Among them is a theorem
whose statement is \lean{False}, in a module that compiles and for which \lean{#print axioms}
reports that the theorem depends on no axioms. A successful build and a clean list of axioms
do not, then, by themselves show that a theorem holds; they show it for code that does not
write into the environment behind the kernel's back and that imports only compiled files
that Lean itself wrote. A replay covers the rest, within
what a replay examines. Our program rejects all \hostilebad{}, each with failures recorded
under exactly the names its design predicts, and accepts the \hostilecontrols{} controls.
Lean's checker with
\lean{--fresh} rejects \hostilefreshrejects{} of the \hostilebad{}, and in its ordinary mode,
which examines the module it is given and trusts that module's imports,
\hostiledefaultrejects{}. In the table a checker \emph{rejects} when it exits with an error.

% Generated by tools/numbers.py from verify/results/hostile.txt and independent-hostile.txt. Do not edit by hand.
\begin{table}[tp]
\centering\small
\begin{tabular}{@{}p{0.5\textwidth}cccc@{}}
\toprule
The compiled file holds & \multicolumn{2}{c}{\lean{leanchecker}} & \lean{lean4lean} & ours\\
 & ordinary & \lean{--fresh} & \lean{--fresh} & \\
\midrule
nothing wrong: a theorem, a definition, an inductive type, a structure & accepts & accepts & accepts & accepts\\
nothing wrong: the tool the other modules use to write into a compiled file & accepts & accepts & accepts & accepts\\
a proof of \lean{False} from a declared axiom, which the list of axioms shows & accepts & accepts & accepts & accepts\\
a proof by \lean{sorry}, which the list of axioms shows & accepts & accepts & accepts & accepts\\
a theorem \lean{False} whose proof is \lean{True.intro}, added with no check & \textbf{rejects} & \textbf{rejects} & \textbf{rejects} & \textbf{rejects}\\
a definition whose value is of another type, added with no check & \textbf{rejects} & \textbf{rejects} & \textbf{rejects} & \textbf{rejects}\\
a theorem whose proof uses an unsafe definition & \textbf{rejects} & \textbf{rejects} & \textbf{rejects} & \textbf{rejects}\\
a recursor whose two rules were exchanged after the kernel made it & \textbf{rejects} & \textbf{rejects} & \textbf{rejects} & \textbf{rejects}\\
a constructor given a false number of fields & \textbf{rejects} & \textbf{rejects} & \textbf{rejects} & \textbf{rejects}\\
a theorem two modules declare; the version Lean keeps has no proof & accepts & \textbf{rejects} & \textbf{rejects} & \textbf{rejects}\\
the same; the version Lean sets aside has no proof & accepts & accepts & accepts & \textbf{rejects}\\
an inductive type that lists a constructor that does not exist & \textbf{rejects} & \textbf{rejects} & \textbf{rejects} & \textbf{rejects}\\
a definition stored under the name of another constant & \textbf{rejects} & \textbf{rejects} & \textbf{rejects} & \textbf{rejects}\\
an inductive type with no recursor, and a theorem that uses one & accepts & accepts & accepts & \textbf{rejects}\\
an inductive type whose record was altered after the kernel made it & accepts & accepts & accepts & \textbf{rejects}\\
a constant of the quotient kind with the type \lean{False} & accepts & accepts & accepts & \textbf{rejects}\\
a theorem \lean{False} proved from that constant, with no axioms & accepts & \textbf{rejects} & \textbf{rejects} & \textbf{rejects}\\
a constant of the inductive kind with the type \lean{False} & accepts & accepts & accepts & \textbf{rejects}\\
an inductive type whose listed constructor is an unsafe definition & accepts & accepts & accepts & \textbf{rejects}\\
a theorem \lean{False} in the first part of a module-system file only & accepts & accepts & accepts & \textbf{rejects}\\
the same, the part replayed marked as outside the module system & accepts & accepts & accepts & \textbf{rejects}\\
in the first part, an axiom \lean{False} under the name of a true theorem & accepts & accepts & accepts & \textbf{rejects}\\
in the first part, an axiom with other universe parameters & accepts & accepts & accepts & \textbf{rejects}\\
in the first part, an axiom \lean{False} for a partial definition, not marked unsafe & accepts & accepts & accepts & \textbf{rejects}\\
the same for an unsafe definition & accepts & accepts & accepts & \textbf{rejects}\\
the same for an unsafe opaque constant & accepts & accepts & accepts & \textbf{rejects}\\
a first part whose import list names one module more & accepts & accepts & accepts & \textbf{rejects}\\
a theorem \lean{False} in the file of compiled code & accepts & accepts & accepts & \textbf{rejects}\\
a file of compiled code whose import list names one module more & accepts & accepts & accepts & \textbf{rejects}\\
\bottomrule
\end{tabular}
\caption{The modules built to test the checkers, and what each checker does with each. The first \hostilecontrols{} are controls, which a checker should accept; a sound replay rejects the others. The ordinary mode of \texttt{leanchecker} examines the named module and trusts its imports; \texttt{lean4lean} is the independent kernel of \cref{sec:checker}.}\label{tab:hostile}
\end{table}


\emph{Agreement with Lean's checker.} Besides the closure of \lean{Main},
\lean{leanchecker --fresh} was run on
the module of Mathlib that declares the statement \lean{RiemannHypothesis}, and on Mathlib as
a whole; it accepted both. Our program
accepts the same closures; those two have \stockreplayrzdecls{} and
\stockreplaymathlibdecls{} declarations. All of this ran on the second machine's build.
\ifsecondrun
In its ordinary mode \lean{leanchecker} takes a name as a prefix: it checks the module of
that name and every module whose name extends it, each module's own declarations against
its imports, importing the closure again for each and all at once. We ran it on
\stocksample{} names, one process a name and \stocksampleatatime{} at a time: \lean{Main}
and \stocksampledrawn{} taken from the modules that hold a declaration of the cone by a
fixed selection that can be repeated (\lean{shuf} with a constant in place of random bytes;
the names are listed in the results). \stockprefixnote{} \stocksamplenote{} The checker also
accepts the module \lean{Statements}, which holds the archive's own definitions of the
statements and is not among the modules a sample could draw.
\fi

\emph{The files of the module system.} A module that uses Lean's module system is compiled to
five files: a first part, a part for the language server, a private part, and two files of
compiled code. A file outside the system, such as every file of the archive, reads the
private part, and that is what a replay replays; a module inside the system that imports
with \lean{import all} reads it too. Otherwise a module inside the system receives its
constants from the first part, from the modules that the first part's import list names,
and, when a module is loaded only for its compiled code, from a file of compiled code and
from the modules that its import list names. Lean writes the first part from the
declarations of the private part, each unchanged, reduced to an axiom with the same
statement, or left out; it gives the first part and that file of compiled code the import
list of the private part; and the command \lean{lean}, which compiles a file, writes no
constants into the files of compiled code. Our program fails unless all of that holds. In Mathlib's
closure every one of the \stockexportedmodules{} modules uses the module system, and the
\stockexported{} constants of their first parts are what Lean derives from the parts
replayed, with one qualification. For a recursive definition that is compiled to code Lean
stores a second constant, marked partial, which a replay does not examine and which no
theorem outside the module system may use. Lean \leanversion{} writes its statement into the
first part as an axiom marked unsafe; in the files of Lean's own library that Mathlib
imports, as distributed,
\stockpartialunmarked{} such axioms are not marked, so that a theorem of a module inside the
system may use one, and \lean{#print axioms} would not list it. The program allows the
unmarked form in those files and nowhere else, and counts it. This comparison rests on the
fourth review's reading of how Lean \leanversion{} imports; it is not a claim that no other
route exists.

\emph{A second implementation of the kernel.} All of the above uses the kernel inside Lean.
\lean{lean4lean} \cite{lean4lean} is an implementation of the type checker written in Lean.
Its authors describe it as derived directly from the C++ code of the kernel and as likely to
share errors with it, so it is a second implementation and not a fully independent one. We
used a fixed commit of it, made for a release candidate of Lean 4.33.0, and built it with
Lean \leanversion{}. In its source the terms are Lean's own, with Lean's functions for
substitution, and compiled files are read by Lean's reader (\lean{readModuleDataParts} in
its \lean{Replay.lean}); type inference and definitional equality
(\lean{TypeChecker.lean}), the checking of inductive types (\lean{Inductive/Add.lean}) and
the quotient type (\lean{Quot.lean}) are written there in Lean. We ran it where its speed
allowed, on the second machine. On the \hostilebad{} modules that a replay must reject, with
\lean{--fresh}, it
rejects \independentfreshrejects{}; like Lean's checker, it does not compare the inductive
types and the other constants it builds with those of the files (\cref{tab:hostile}). With
\lean{--fresh} on the module of Mathlib that declares \lean{RiemannHypothesis} and
\lean{riemannZeta}, and on the one that declares the $L$-function of a Dirichlet character,
it accepts the whole import closure of each, in that machine's build of Mathlib: every
definition on which those three depend was accepted by a second implementation of the
kernel.
\ifsecondrun
It was also run on the same \independentsample{} names, one process a name and
\independentsampleatatime{} at a time; it takes a name in the same way.
\independentsamplenote{} It accepts the module \lean{Statements} too. It was not run on the
whole archive.
\else
It has not been run on the archive.
\fi

\subsection{What this establishes, and what it does not}
In Lean \leanversion{} with Mathlib at the pinned commit, \cref{st:rh,st:grh,st:strip} are
theorems\ifsecondrun, with no assumption beyond the standard axioms\else; by the lists of
axioms recorded in the compiled files they assume nothing beyond the standard axioms, and
the replay establishes what \cref{sec:replay} says\fi. That is the whole of what the machine
certifies. It does not certify any description of the proof, including ours.

\ifsecondrun
Three limits should be stated. First, the full replay uses one kernel, the one in Lean
\leanversion{}: the build, our replay and Lean's checker share it, together with the code
that reads compiled files, so their agreement is one witness to the kernel's judgment and not
several. That code takes on trust the hashes and flags that names and terms carry in a
compiled file; nothing recomputes them, and the third review describes how a file written
by hand could use that to make a statement refer to another constant than the one its name
shows. The files replayed here were written by Lean in our second build or, for Lean's own
library, are those of the published release; and \lean{check.sh} writes out
\cref{st:rh,st:grh,st:strip} and elaborates each afresh against its theorem. The second
implementation covered the definitions named above and a sample, not the whole.
\else
Four limits should be stated. First, the replay of the whole archive reported here was made
with the earlier form of our program, and establishes what \cref{sec:replay} says and no
more; and the axioms are those of the lists recorded in the compiled files
(\cref{sec:axioms}), which for Mathlib were made by whoever built the cache. The
corrected program's run, which recomputes both, is under way. Second, the full replay uses
one kernel, the
one in Lean \leanversion{}: the build, our replay and Lean's checker share it, together with
the code that reads compiled files, so their agreement is one witness to the kernel's
judgment and not several. The second kernel covered the statement, not the archive.
\fi
The kernel
inside Lean also takes on trust the definitions of a few constants of Lean's own library that
it recognises by name, among them the arithmetic of the natural numbers. On the second
machine we compared the toolchain's files with the published release, and the second
implementation ran on closures in which those definitions are declared; when it adds one
of them it checks that the definition has the form assumed (\lean{checkPrimitiveDef} in its
source). Neither was done on the first machine. A run of another implementation of the
kernel over the whole archive is the check most worth adding.
\ifsecondrun
Second, the replay program has been reviewed and tested as
described, but by the model that wrote it and by no person. That the kernel accepts each of
the \replaydecls{} constants no longer rests on it, since Lean's own checker finds the same;
the comparisons and the further versions listed in \cref{sec:checker} do. Third, both builds are ours; we
know of no other at the time of writing.
\else
Third, the replay program has been reviewed and tested as described, but by the model that
wrote it and by no person. Fourth, the build is ours; we know of no other at the time of
writing.
\fi
